\begin{problem}[From a sheaf on a basis to a sheaf on the whole space]
Let $\mathcal B$ be a basis of $X$ closed under finite intersections (or replace intersections by basis refinements). Suppose $F$ assigns groups and restrictions to members of $\mathcal B$ and satisfies the sheaf axiom for covers by basis elements. Show that there is a sheaf $\widetilde F$ on $X$, unique up to unique isomorphism, whose restriction to the basis is $F$.
\end{problem}

\paragraph{Why this problem matters.}
Algebraic geometry usually specifies regular functions first on a convenient basis, not on every open set.

\paragraph{Useful ideas before starting.}
Use the universal property, stalk preservation, and the fact that a morphism of sheaves is an isomorphism when it is so on every stalk.  Keep the distinction between a presheaf statement on sections and a local statement on germs explicit.

\paragraph{Guided hints.}
For arbitrary $U$, define sections as compatible families on all basis opens $B\subseteq U$, or first form the right Kan-style presheaf and sheafify.

\ifdefined\FullProblemDossiers
\begin{solution}
Define
\[
\widetilde F(U)=\left\{(s_B)_{B\in\mathcal B,\,B\subseteq U}: s_B\in F(B),\; s_B|_C=s_C\text{ whenever }C\subseteq B\right\}.
\]
Restrictions forget components. Compatible families on an open cover glue componentwise: for a basis element $B\subseteq U$, the sets $B\cap U_i$ are covered by basis elements, and the basis sheaf axiom glues the corresponding local data uniquely to $s_B$. These $s_B$ are mutually compatible, so $\widetilde F$ is a sheaf. If $U=B$ is itself in the basis, the map $F(B)\to\widetilde F(B)$ sending a section to all its restrictions is an isomorphism by the basis sheaf axiom. Any other extension has the same sections on arbitrary opens because a sheaf is determined by compatible sections on a basis cover. Thus the extension is unique up to unique isomorphism.
\end{solution}

\paragraph{What happened mathematically.}
A sufficiently rich basis contains all local information needed to reconstruct the sheaf.

\paragraph{Alternative approaches.}
Whenever possible one may replace the espace-etale model by the two-plus construction, or conversely replace a cover/refinement argument by the universal property.  The two approaches agree because both construct the same associated sheaf.

\paragraph{Extensions.}
Apply this to $D(f)\mapsto A_f$ in VI/17.

\paragraph{Related problems.}
Compare with VI.14.P1 for the complete legacy construction and with the later sheaf chapters for operations that are deliberately not migrated here.

\paragraph{Provenance.}
New synthesis problem preparing the structure-sheaf construction without consuming VI/17.
\else
\paragraph{Provenance.}
New synthesis problem preparing the structure-sheaf construction without consuming VI/17.
\fi
